\documentclass[11pt,a4paper]{amsart}

\usepackage[margin=1in]{geometry}
\usepackage{amsmath,amssymb,amsthm,mathtools}
\usepackage[expansion=false]{microtype}
\usepackage[colorlinks=true,linkcolor=black,citecolor=black,urlcolor=black]{hyperref}
\usepackage[nameinlink,capitalize]{cleveref}

\hypersetup{
  pdftitle={Class recovery and inflation for repeated-block skew diagrams},
  pdfauthor={John Shields},
  pdfkeywords={skew Schur function, skew-equivalence, repeated-block diagram,
    ribbon Schur function, profile inflation, marked composition, class recovery}
}
\pdfinfo{ /ORCID (https://orcid.org/0009-0004-3882-4571) }

\numberwithin{equation}{section}
\theoremstyle{plain}
\newtheorem{theorem}{Theorem}[section]
\newtheorem{proposition}[theorem]{Proposition}
\newtheorem{lemma}[theorem]{Lemma}
\newtheorem{corollary}[theorem]{Corollary}
\theoremstyle{definition}
\newtheorem{definition}[theorem]{Definition}
\newtheorem{remark}[theorem]{Remark}
\newtheorem{question}[theorem]{Question}

\title[Repeated-block classes and inflation]
{Class recovery and inflation\\
for repeated-block skew diagrams}
\author{John Shields}
\date{October 4, 2026}
\thanks{Independent researcher, Ginza 2-15-2, KR Ginza II 5F, Chuo-ku, Tokyo, Japan 104-0061.\\
\hspace*{\parindent}Corresponding author: \href{mailto:math@johnnyshields.com}{math@johnnyshields.com}.  ORCID: \href{https://orcid.org/0009-0004-3882-4571}{0009-0004-3882-4571}.\\
\hspace*{\parindent}This work is archived at \href{https://doi.org/10.5281/zenodo.23133727}{doi:10.5281/zenodo.23133727}.}
\subjclass[2020]{Primary 05E05; Secondary 05A17, 15A15}
\keywords{skew Schur function, skew-equivalence, repeated-block diagram,
ribbon Schur function, profile inflation, marked composition, class recovery}

\begin{document}

\begin{abstract}
We recover complete connected skew Schur classes for two specified repeated-block constructions.  For separated constant-gap inners repeated at least twice, every mate comes from the outer ribbon class and global rotation, without assuming that the inner diagram is rigid.  Symmetric inners remove the otherwise separate orientation copy.  For an asymmetric two-row inner, adding the same nonnegative integer to every row length and adjacent overlap-minus-one preserves the full class below the shorter row's excess over the inner gap.  At that exact threshold, and at every larger inflation, only rotation remains.  Inflation is injective: the images of formerly equal diagrams split into distinct classes.  Total contraction proves the first inverse; a lowest-degree boundary determinant with private column residues proves the second.  Both theorems quantify over arbitrary connected mates.
\end{abstract}

\maketitle

\section{Introduction}\label{sec:introduction}

A repeated context can determine its inner blocks even when their individual
skew Schur functions do not.  We prove this for a separated constant-gap
construction, recovering every connected mate from an outer ribbon class.
For a different construction with a two-row inner, we determine exactly how
long that outer ambiguity survives additive inflation of the profile.

Diagrams are identified up to translation.  Write \(D^*\) for rotation
through \(180\) degrees and
\[
 \operatorname{Cl}(D)=\{H:H\text{ is connected and }s_H=s_D\}.
\]
A profile \((r_1,\ldots,r_n;g_1,\ldots,g_{n-1})\) records row lengths
from top to bottom and adjacent overlaps minus one.  The join
\(P\cdot_d Q\) concatenates the row words with joining gap \(d\).
The merger \(P\odot_d Q\) replaces the two adjoining endpoint rows
of lengths \(u,v\) by one row of length \(u+v-d\), retaining the
other rows and gaps.  These are profile operations, distinct from marked
composition.  They are used only when their profiles define connected
skew diagrams.  A superscript on a number or a word inside a profile
denotes repetition.

For a positive composition \(\alpha\), let \(r_\alpha\) be its ribbon
Schur function, indexing the ribbon's rows from top to bottom, and put
\[
 \mathcal R(\alpha)=\{\beta:r_\beta=r_\alpha\}.
\]
The classification of these classes by reversal of composition factors
is stated in \cite[Cor.~4.2]{Billera2006DecomposableCompositions},
with the necessity direction supplied by
\cite{Shields2026Coarsenings}.
Both papers read ribbon rows from bottom to top.  Our convention reverses
those diagrams and leaves \(r_\alpha\) unchanged; all literal marked
presentations below use the top-to-bottom convention.

Fix integers \(k\ge2\) and \(0\le c<d\), and a constant-gap inner
\begin{equation}\label{eq:saturation-domain}
 \begin{gathered}
 A=(a,I,b;c^{k-1}),\\
 a,b>c+d,\qquad
 I=(a_2,\ldots,a_{k-1}),\qquad
 a_i>\max(2c,d)\quad(2\le i\le k-1).
 \end{gathered}
\end{equation}
The word \(I\) is empty when \(k=2\).  Set \(z=a+b-d\) and
\begin{equation}\label{eq:repeated-blocks}
 A_m=(a,I,(z,I)^{m-1},b;c^{m(k-1)}),\qquad
 F_\alpha=A_{\alpha_1}\cdot_d\cdots\cdot_d A_{\alpha_\ell}.
\end{equation}
Thus \(A_m=A^{\odot_d m}\).

\begin{theorem}[Repeated-block saturation]\label{thm:repeated-block-saturation}
Under \eqref{eq:saturation-domain}, let \(\alpha\) be a positive
composition with \(M=|\alpha|\ge2\).  Then
\begin{equation}\label{eq:saturation-class}
 \operatorname{Cl}(F_\alpha)
 =\{F_\beta,F_\beta^*: \beta\in\mathcal R(\alpha)\}.
\end{equation}
Its cardinality is
\begin{equation}\label{eq:saturation-cardinality}
 |\operatorname{Cl}(F_\alpha)|=
 \begin{cases}
 |\mathcal R(\alpha)|,&A=A^*,\\
 2|\mathcal R(\alpha)|,&A\ne A^*.
 \end{cases}
\end{equation}
After at most one global rotation, every member has the literal
case-(a) marked presentation \(\beta\circ_{(d)}A\), with
\(\beta\in\mathcal R(\alpha)\) and both end-row marks of length
\(d\).  These charts satisfy Hypotheses I--IV of
\cite{McNamara2009Classification}.  No rigidity assumption is imposed on
\(A\).
\end{theorem}

Cutting the larger seams recovers the component functions and their row
counts.  It does not recover a component consisting of one nonrigid copy
of \(A\).  Contracting all those seams instead gives \(A_M\), which
is rigid for \(M\ge2\).  The recovered component sizes locate the cuts
in its ordered row word, and the unordered endpoint pairs determine their
splits.  This recovers even the one-copy components literally.

The inflation theorem has a different domain.  Fix integers
\(1\le c<a<b\), and put
\begin{equation}\label{eq:thin-domain}
 \begin{gathered}
 p=a-c,\qquad q=a+b-c,\\
 E_m=(a,(a+b)^{m-1},b;c^m),\qquad
 G_\alpha=E_{\alpha_1}\cdot_0\cdots\cdot_0 E_{\alpha_\ell}.
 \end{gathered}
\end{equation}
For an integer \(t\ge0\), define profile inflation by
\begin{equation}\label{eq:inflation}
 I_t(r;g)=(r_1+t,\ldots,r_n+t;g_1+t,\ldots,g_{n-1}+t).
\end{equation}

\begin{theorem}[Sharp inflation threshold]\label{thm:sharp-inflation}
Under \eqref{eq:thin-domain}, for every nonempty composition \(\alpha\)
and every integer \(t\ge0\),
\begin{equation}\label{eq:inflation-class}
 \operatorname{Cl}(I_tG_\alpha)=
 \begin{cases}
 I_t\bigl(\operatorname{Cl}(G_\alpha)\bigr)
 =\{I_tG_\beta,(I_tG_\beta)^*: \beta\in\mathcal R(\alpha)\},
       &0\le t<p,\\[2mm]
 \{I_tG_\alpha,(I_tG_\alpha)^*\},&t\ge p.
 \end{cases}
\end{equation}
The respective cardinalities are \(2|\mathcal R(\alpha)|\) and two.
For \(0<t<p\), the first class has legal presentations
\(\beta\circ_{(t)}E(t)\), where \(E(t)=(a+t,b+t;c+t)\).
At \(t=0\), use the horizontal empty-overlap convention.
For \(0\le t<p\), ribbon factor reversals and global rotation generate
every member of the first class
\cite{Shields2026Coarsenings}.
\end{theorem}

Inflation is injective on profiles and commutes with rotation.  At the
threshold, it does not identify previously distinct diagrams: their images
cease to have equal Schur functions and split into rotation classes.
There is an actual decrease in class size only when
\(|\mathcal R(\alpha)|>1\).  Neither main theorem is a special case
of the other, and no commutation statement between their constructions
is asserted.
General inflation gives only one-sided no-escape; erosion of a connected
zero-free nonline diagram instead gives an exact class bijection, with
legal charts transported by the alternatives in \Cref{lem:inflate-chart}.

The two main theorems give the separated repeated-block inverse without
inner rigidity and the exact threshold for additive profile inflation.
The latter is not the multiplication of beta
numbers studied in the cofactor companion \cite{Shields2026Cofactor}.
The overlap invariants used below belong to Reiner, Shaw and van
Willigenburg \cite{Reiner2007Coincidences}.  The Toeplitz-minor
correspondence and its row/column reversal interpretation are due to
Garc\'ia-Garc\'ia and Tierz \cite{GarciaGarcia2020ToeplitzMinors}.

\Cref{sec:preliminaries} recalls the coefficient readings and fixes the
composition conventions.  \Cref{sec:profile-transport} proves erosion,
filtered lowering and exact transport on zero-free classes.
\Cref{sec:overlap-powers} proves rigidity of the total contractions, and
\Cref{sec:saturation-proof} reconstructs the literal blocks.
\Cref{sec:thin-and-lowering} gives the two-row class inverse.
\Cref{sec:inflation-proof} recovers outer cuts at the first boundary and
proves the threshold theorem.
\Cref{sec:examples} records examples and the scope of the auxiliary checks.

\section{Profiles, cuts, and ribbon substitution}\label{sec:preliminaries}

Work in \(\Lambda=\mathbb Q[h_1,h_2,\ldots]\), with \(h_0=1\)
and \(h_j=0\) for \(j<0\).  Weighted degree assigns \(h_j\) degree
\(j\); ordinary degree assigns every positive-index generator degree one.
For a diagram with \(n\) nonempty rows, normalize its strict beta lists
\[
 X=(x_1>\cdots>x_n),\qquad Y=(y_1>\cdots>y_n=0).
\]
Then
\begin{equation}\label{eq:beta-profile}
 s_D=\det(h_{x_i-y_j}),\qquad
 r_i=x_i-y_i,\qquad g_i=x_{i+1}-y_i.
\end{equation}
The exact profile inequalities are
\begin{equation}\label{eq:profile-validity}
 r_i>0,\qquad 0\le g_i<\min(r_i,r_{i+1}).
\end{equation}
They characterize connected skew diagrams up to translation.  The lists
are recovered recursively from \(x_n=r_n\), \(y_n=0\), and
\(y_i=x_{i+1}-g_i\), \(x_i=y_i+r_i\).  In particular,
\begin{equation}\label{eq:content-span}
 N_D=x_1=\sum_i r_i-\sum_i g_i.
\end{equation}
All beta lists in this paper are full lists: common entries are not
canceled unless a deletion is explicitly specified.  A diagram is
zero-free when its beta lists are disjoint, so its Jacobi--Trudi matrix
has no unit entry.  Rotation reverses both profile words and preserves
the Schur function.

\subsection{Readings on an arbitrary mate}\label{subsec:readings}

The polynomial-readings lemma of \cite{Shields2026Cofactor} applies to
any connected diagram, before a Hessenberg condition has been established.
It says that \(s_D\) determines the row count as its maximum ordinary
degree, and the row multiset as the indices of its unique minimum-square
monomial, with coefficient one.  Here the square weight of
\(\prod h_{j_i}\) is \(\sum j_i^2\).  It also reads the largest
generator index \(N_D\) and its signed coefficient
\begin{equation}\label{eq:maximal-cofactor}
 (-1)^{n-1}[h_{N_D}]s_D
 =\det(h_{x_{i+1}-y_j})_{i,j=1}^{n-1}\ne0.
\end{equation}
The minimum-square monomial of this cofactor, padded with zeros to length
\(n-1\), reads the gap multiset.  These are polynomial forms of the
row-overlap data of \cite{Reiner2007Coincidences}.

\begin{lemma}\label{lem:hessenberg-membership}
Let \(D\) be a connected diagram with \(n\) rows and beta lists as in
\eqref{eq:beta-profile}.  Its signed maximal-label coefficient is a
monomial if and only if its full Jacobi--Trudi board is upper Hessenberg:
\begin{equation}\label{eq:hessenberg-membership}
 \begin{gathered}
 (-1)^{n-1}[h_{N_D}]s_D\text{ is a monomial}\\
 \Longleftrightarrow\quad
 x_i-y_j<0\text{ for }i\ge j+2
 \quad\Longleftrightarrow\quad
 r_i>g_{i-1}+g_i\quad(2\le i\le n-1).
 \end{gathered}
\end{equation}
\end{lemma}

\begin{proof}
The entries of the full board below its first subdiagonal are the entries
of the cofactor in \eqref{eq:maximal-cofactor} on and below its own first
subdiagonal.  That subdiagonal has indices
\(x_{i+2}-y_i=g_i+g_{i+1}-r_{i+1}\), and indices strictly decrease down
each column, so the last two conditions in \eqref{eq:hessenberg-membership}
agree.  Split the cofactor at the negative entries of its first
subdiagonal; strictness makes the full lower-left rectangles zero.  A zero
diagonal contributes an isolated unit.  Each remaining block of order at
least two has its distinct nonzero diagonal monomial and a term containing
its private largest label, whose coefficient is nonzero by strict
rearrangement.  It cannot be a monomial.  A product of nonzero polynomials
can be a monomial only if every factor is.  Thus all cofactor blocks must
be singleton, precisely the stated inequalities.
\end{proof}

Thus equality with a Hessenberg source forces Hessenberg form, as well
as the same row and gap multisets.  The cofactor in that case is
\(\prod_i h_{g_i}\).  Its omitted zero factors are recovered from the
row count.

A nonempty connected skew Schur polynomial is irreducible in \(\Lambda\).
We use the corrected irreducibility theorem
\cite{Barekat2009CoincidencesCorrigendum}, rather than an assertion that
specialization preserves irreducibility.  Consequently a product of
functions of connected cut components determines their functions with
multiplicity.  The minimum-square monomial, whose coefficient is one,
removes any scalar ambiguity between associated factors.

The other reconstruction input is the theorem on two boundary observations
in \cite{Shields2026SkewSchur}.  If \(D,H\) are connected and
\(\overline D,\overline H\) are obtained by deleting their entire first
rows, then
\begin{equation}\label{eq:boundary-reconstruction-input}
 s_D=s_H,\qquad s_{\overline D}=s_{\overline H}
 \quad\Longrightarrow\quad D=H.
\end{equation}
The conclusion is equality up to translation with the displayed orientation;
\(s_\varnothing=1\).  Its proof is not repeated here.

\subsection{Seam coefficients}\label{subsec:seam-coefficients}

For an upper-Hessenberg profile, the interval expansion of
\cite{Shields2026NegativeCount} is
\begin{equation}\label{eq:interval-expansion}
 s_D=\sum_{T\subseteq\{1,\ldots,n-1\}}
 (-1)^{|T|}\prod_{i\in T}h_{g_i}
 \prod_{J\in\pi(T)}h_{R_J},\qquad
 R_J=\sum_{j\in J}r_j-\sum_{j,j+1\in J}g_j.
\end{equation}
Here \(\pi(T)\) is the interval partition obtained by contracting the
seams indexed by \(T\).  Every interval label \(R_J\) is positive
and exceeds every row length in \(J\) when \(|J|>1\).  The strict
increase follows by extending an interval one row at a time and using
\eqref{eq:profile-validity}.

\begin{lemma}\label{lem:seam-coefficients}
Let \(D\) be connected and upper Hessenberg, and let \(d>0\) be an
integer.  Suppose every row exceeds \(d\) and exactly \(s\) gaps equal
\(d\).  Cut those seams to obtain \(D_1,\ldots,D_{s+1}\), and
contract them all to obtain \(C\).  Then these profiles are connected
and upper Hessenberg, and
\begin{equation}\label{eq:seam-extremes}
 [h_d^0]s_D=\prod_j s_{D_j},\qquad
 [h_d^s]s_D=(-1)^s s_C.
\end{equation}
None of the polynomials on the right involves \(h_d\).

For an arbitrary connected upper-Hessenberg \(D\), set
\begin{equation}\label{eq:unit-homogenization}
 \widehat s_D(u)=u^n s_D(h_1/u,h_2/u,\ldots).
\end{equation}
If \(s\) gaps are zero, cutting or contracting all zero seams gives
\begin{equation}\label{eq:unit-extremes}
 [u^0]\widehat s_D(u)=\prod_j s_{D_j},\qquad
 [u^s]\widehat s_D(u)=(-1)^s s_C.
\end{equation}
Here the cut components and the total contraction have no zero gaps.
\end{lemma}

\begin{proof}
Every interval label exceeds \(d\), so \(h_d\) occurs in
\eqref{eq:interval-expansion} only at the specified seams.  Its constant
and highest coefficients respectively cut and contract all those seams,
giving \eqref{eq:seam-extremes}.  This is the safe-cut argument of
\cite{Shields2026NegativeCount}; only separation from the interval labels
is needed, and other gaps may lie on either side of \(d\).

Cutting removes an incident gap.  At a contraction with endpoint rows
\(v,w\) and exterior gaps \(e,f\), the inequalities
\(v>e+d\), \(w>d+f\) give \(v+w-d>e+f\); an absent exterior
gap is zero.  Thus the contracted profile remains connected and
Hessenberg.  Its new row exceeds \(d\), so no new \(h_d\) occurs.

Homogenization in \eqref{eq:unit-homogenization} replaces \(h_0\)
by \(u\), leaving positive entries unchanged and negative entries zero.
The only units on a Hessenberg board are its zero-gap subdiagonal
entries.  The same seam selection, with these units retained, gives
\eqref{eq:unit-extremes}; the contraction calculation also applies at \(d=0\).
\end{proof}

For later use, the cut--merge identity on this domain is
\begin{equation}\label{eq:cut-merge}
 s_{P\cdot_d Q}=s_Ps_Q-h_ds_{P\odot_d Q}.
\end{equation}
When the incident interior gaps are \(e,f\), the positive-gap version
requires endpoint lengths \(v>e+d\) and \(w>d+f\), in addition to
Hessenberg form inside the two pieces.  Connectedness of the join alone
is not sufficient.  Equation~\eqref{eq:cut-merge} is the corresponding
case of the same interval expansion.

\subsection{Ribbon substitution}\label{subsec:ribbon-substitution}

If \(\gamma\) is obtained from \(\alpha\) by merging adjacent
parts, write \(\gamma\succeq\alpha\).  The classical ribbon formula is
\begin{equation}\label{eq:ribbon-coarsening}
 r_\alpha=\sum_{\gamma\succeq\alpha}
 (-1)^{\ell(\alpha)-\ell(\gamma)}
 h_{\gamma_1}\cdots h_{\gamma_{\ell(\gamma)}}.
\end{equation}
We use it together with the factor-reversal classification in
\cite{Shields2026Coarsenings}.  For compositions, ordinary
near-concatenation \(\rho\odot\sigma\) merges the last and first
parts by addition, and
\begin{equation}\label{eq:ribbon-composition}
 \gamma\circ\rho
 =\rho^{\odot\gamma_1}\cdot\ldots\cdot
   \rho^{\odot\gamma_{\ell(\gamma)}}.
\end{equation}
Two ribbon functions are equal precisely when their compositions admit
factorizations whose corresponding factors agree up to reversal.
In particular, equality preserves the number of parts and their multiset;
this also follows from the maximum ordinary-degree term in
\eqref{eq:ribbon-coarsening}.

\begin{lemma}\label{lem:triangular-substitution}
Let \(K\) be a field, and let \(B_m\) be polynomials in independent
variables \(z_1,z_2,\ldots\) over \(K\).  Suppose
\[
 B_m=c_m z_{N_m}+d_m,\qquad
 0<N_1<N_2<\cdots,
\]
where \(c_m,d_m\in K[z_1,\ldots,z_{N_m-1}]\) with \(c_m\ne0\); the
case \(d_m=0\) is allowed.  Then substitution \(H_m\mapsto B_m\) defines an
injective homomorphism \(K[H_1,H_2,\ldots]\to K[z_1,z_2,\ldots]\).
\end{lemma}

\begin{proof}
Take a nonzero relation with largest index \(m\), and expand it as a
polynomial in \(B_m\).  By induction, its leading coefficient, a
nonzero polynomial in \(B_1,\ldots,B_{m-1}\), is nonzero.  The highest
power of \(z_{N_m}\) has that coefficient multiplied by a power of
\(c_m\), and cannot vanish in the polynomial domain.
\end{proof}

For a constant-gap Hessenberg profile \(D=(r_1,\ldots,r_n;c^{n-1})\),
put \(\rho=(r_1-c,\ldots,r_n-c)\).  Equation~\eqref{eq:interval-expansion}
gives
\begin{equation}\label{eq:constant-gap-substitution}
 s_D=h_c^n\,
 r_\rho\left(\frac{h_{c+1}}{h_c},
             \frac{h_{c+2}}{h_c},\ldots\right).
\end{equation}
For \(c>0\), this is read in the localization at \(h_c\), and its
substitution is injective because the generators \(h_{c+j}\) are
independent.  For \(c=0\), it is the ordinary ribbon formula.
This faithful substitution compares diagrams already known to have
constant gap.  It does not by itself expose an arbitrary mate.

\subsection{Marked geometry}\label{subsec:marked-geometry}

We use the full marked geometry of \cite{Shields2026Gluing}, following
McNamara and van Willigenburg \cite{McNamara2009Classification}.
A marked inner \(\omega\) consists of a connected diagram \(E\),
two congruent connected end copies \(W_-=W_{\mathrm{sw}}\) and
\(W_+=W_{\mathrm{ne}}=W_-+v\), and their attachment case.
The lower copy contains the southwest cell of \(E\), and the upper
copy contains its northeast cell.  Put
\[
 O=E\setminus(W_-\cup W_+),\qquad w=N_W,\qquad q=N_E-w=c(v),
 \qquad e=(1,1),
\]
where \(c(i,j)=j-i\), rows increase downward and columns to the right.
The four cases, in lower/upper order, are
\[
 \begin{array}{c|cccc}
 \text{case}&(a)&(b)&(c)&(d)\\ \hline
 \text{attachments}&(H,H)&(V,V)&(H,V)&(V,H).
 \end{array}
\]
A horizontal lower attachment uses the western neighbor of the
southwest middle cell, and a vertical one its southern neighbor.
The upper attachment uses the eastern or northern neighbor of the
northeast middle cell.  Cases \((a),(b)\) are ordinary and
\((c),(d)\) are mixed.

The mark is legal under Hypotheses I--IV:
\begin{enumerate}
\item[(I)] No strictly larger skew diagram containing \(W_-\), on the
 same content diagonals, has translated copies at both ends of \(E\).
\item[(II)] An occupied diagonal meets neither marked copy: \(0\le w<q\).
\item[(III)] Removing either marked copy leaves a connected skew diagram.
\item[(IV)] In \(\boldsymbol E=\bigcup_{k\in\mathbb Z}(E+kv)\), no
 hatted-middle copy is adjacent to a hatted-mark copy.  For copies
 \(O_1,W_1\), these are
 \[
 \widehat O_1=\{x\in O_1:x+e\in O_1\},\qquad
 \widehat W_1=\{x\in\boldsymbol E:x+e\in W_1\}
       \cup\{x\in W_1:x+e\in\boldsymbol E\}.
 \]
\end{enumerate}
Adjacency means sharing an edge.  The notation \(\overline O,\overline W\)
for these strips is also used in the repeated-block arguments below.
At \(w=0\) the marks are the west edge of the southwest cell and the east
edge of the northeast cell, with horizontal displacement
\[
 v_H=\operatorname{ne}(E)-\operatorname{sw}(E)+(0,1).
\]
The vertical edge reading has \(v_V=v_H-e\).  Conversion to the
horizontal reading antitransposes the outer:
\(D\circ_{\varnothing,V}E=\operatorname{AT}(D)\circ_{\varnothing,H}E\),
where \(\operatorname{AT}(i,j)=(-j,-i)\).
More generally, when both ordinary marks exist, \(v_a-v_b=e\) and
\(D\circ_{\omega_a}E=\operatorname{AT}(D)\circ_{\omega_b}E\).
The usual symmetric-function involution acts on the outer function at
the same time.  A proper chart \(D\circ_\omega E\) has connected factors
with \(|D|,|E|\ge2\).  Rotation transports both marks and their attachments;
it is not an arbitrary Schur-equal replacement of the inner.

For \(d=(i,i+j)\in D\), write \([\ell_j^D,h_j^D]\) for the
content-\(j\) row interval of the outer.  The copy of \(E\) indexed by
\(d\) has origin
\[
 L_\omega(d)=
 \begin{cases}
 jv+ie,&(a),\\
 jv-ie,&(b),\\
 jv+(i-\ell_j^D)e,&(c),\\
 jv-(h_j^D-i)e,&(d).
 \end{cases}
\]
In the mixed cases, a vertical pair on the same northwest, respectively
southeast, boundary ribbon contributes
\(W_++L_\omega(d)+e\), respectively \(W_++L_\omega(d)-e\),
with \(d\) the lower cell.  The ribbon ranks are
\(i-\ell_j^D+1\) and \(h_j^D-i+1\).  These supplemental copies belong
to the actual cell union.

Here are the local geometry and coding facts from
\cite[local geometry and coding lemmas]{Shields2026Gluing} used in
chart transport.  Normalize the southwest cell to \((0,0)\).
For the interval \([\ell_k,h_k]\) on diagonal \(k\), set
\[
 A_k=\ell_{k-1}-\ell_k,\qquad B_k=h_{k-1}-h_k,\qquad
 H_k=h_k-\ell_k+1.
\]
The internal steps \(1\le k<N_E\) are binary and
\(H_k-H_{k-1}=A_k-B_k\).  Conversely positive interval heights,
binary endpoint steps and singleton end diagonals characterize a
connected skew diagram.  In the beta normalization
\eqref{eq:beta-profile}, the full membership words are
\(A_E(j)=\mathbf1_{j\in Y}\), \(B_E(j)=\mathbf1_{j\in X}\);
they agree with the internal steps and are defined for every integer.
Their heights are
\(H_E(j)=\#\{y\in Y:y\le j\}-\#\{x\in X:x\le j\}\).

For a legal positive mark, \(W_-=E\cap(E-v)\).  The middle is thin
on the marked diagonals, and the private diagonals \(w,q-1\) are
singletons.  Its content interval is \([\sigma,q+\tau-1]\), where
\[
 1\le\sigma\le w,\qquad 0\le\tau<w,
 \qquad \tau<\sigma\quad\text{in the mixed cases}.
\]
Thus \(\tau,\sigma\) are the upper and lower middle phases, not
inflation parameters.  In case \((d)\) both marks are
southeast-justified, and in case \((c)\) northwest-justified.
If \(N_D=M\), the composite has span \(N=Mq+w\); its first and
last \(q\) diagonals are translated end windows of \(E\), whose
private overlap reconstructs \(E\).

The boundary words are \(q\)-periodic except at the slots shown below,
whose interior values hold for \(1\le j<M\):
\[
 \begin{array}{c|ll}
 &\text{first slot}&\text{second slot}\\ \hline
 (a)&A_{jq+\tau}=A_j^D&B_{jq+\sigma}=B_j^D\\
 (b)&B_{jq+\tau}=1-A_j^D&A_{jq+\sigma}=1-B_j^D\\
 (c)&B_{jq+\tau}=1-A_j^D&B_{jq+\sigma}=B_j^D\\
 (d)&A_{jq+\tau}=A_j^D&A_{jq+\sigma}=1-B_j^D
 \end{array}
\]
The true initial and terminal values are
\[
 \begin{array}{c|l|l}
 &\text{initial}&\text{terminal}\\ \hline
 (a)&A_\tau=1,\ B_\sigma=0&A_{N-w+\tau}=0,\ B_{N-w+\sigma}=1\\
 (b)&B_\tau=0,\ A_\sigma=1&B_{N-w+\tau}=1,\ A_{N-w+\sigma}=0\\
 (c)&B_\tau=B_\sigma=0&B_{N-w+\tau}=B_{N-w+\sigma}=1\\
 (d)&A_\tau=A_\sigma=1&A_{N-w+\tau}=A_{N-w+\sigma}=0.
 \end{array}
\]
Steps at \(0,N\) are omitted in this table.  Case \((d)\) has
\(B_w=1,B_q=0\), and case \((c)\) has \(A_w=0,A_q=1\).
At empty width the only exceptional residue is zero.

In the two repeated-block families of
\eqref{eq:saturation-domain} and \eqref{eq:thin-domain}, every nonempty mark is a
horizontal row with two horizontal attachments, hence case \((a)\).
An east outer step amalgamates marks; a north step places the next lower
mark one position northwest of the preceding upper mark.  Those family
charts also satisfy Hypothesis V: a marked copy has just one cell
adjacent to the middle.  V and the horizontal-row restriction are not
imposed on the general transport results.

\section{Profile transport}\label{sec:profile-transport}

\subsection{Erosion and filtered lowering}

Let \(\mathcal E(S)=\{(i,j)\in S:(i,j+1)\in S\}\), retaining
the original row coordinates before translation.  Empty rows may occur,
and the result may disconnect.  Write \(r(S)\) for the number of
nonempty rows and call a row or a column a line.  One-step inflation is
\(\mathcal I=I_1\) from \eqref{eq:inflation}: it increases every row
length and adjacent overlap-minus-one by one, or sends the full beta
lists \((X,Y)\) to \((X+1,Y)\).

\begin{lemma}[Erosion]\label{lem:erosion}
If \(S\) has \(r\) rows, take the ordinary-degree-\(r\) part of \(s_S\) and substitute
\(h_j\mapsto h_{j-1}\), including \(h_1\mapsto1\).  The result is
\(s_{\mathcal E(S)}\).  Thus erosion preserves whole Schur equality, and
\(|\mathcal E(S)|=|S|-r\).
\end{lemma}
\begin{proof}
The selected degree keeps precisely permutations using no old unit entry.
After the shift their entries are the physical entries of \((X-1,Y)\).
An old unit becomes a negative zero, and a newly produced index zero remains
one.  This is exactly the determinant of the eroded diagram, including empty
rows.  All other permutations were excluded by their ordinary degree.
\end{proof}

This is a maximum ordinary-degree projection.  It is not a
minimum-degree core extraction, and its substitution \(h_1\mapsto1\)
retains newly empty rows.  The following lowering map instead filters
both the number of positive factors and every positive index.

\begin{lemma}\label{lem:inflation-lowering}
For a connected \(n\)-row diagram \(D\) and an integer \(t\ge1\),
define a linear map on \(\Lambda\) by
\begin{equation}\label{eq:lowering-map}
 \mathcal L_{t,n}\!\left(\sum_\lambda c_\lambda h_\lambda\right)
 =\sum_{\substack{\ell(\lambda)=n\\ \lambda_i\ge t\text{ for all }i}}
 c_\lambda\prod_{i=1}^n h_{\lambda_i-t}.
\end{equation}
Then
\begin{equation}\label{eq:lowering-identity}
 \mathcal L_{t,n}(s_{I_tD})=s_D.
\end{equation}
\end{lemma}

\begin{proof}
Inflation shifts every Jacobi--Trudi index by \(t\).  A term with
all original indices nonnegative then has \(n\) positive factors of
indices at least \(t\), and lowering recovers it, including \(h_0=1\).
If an original negative index becomes positive, it lies in
\(\{1,\ldots,t-1\}\) and fails the index filter; if it becomes a unit,
the term has fewer than \(n\) positive factors.  An index remaining
negative contributes zero.  Thus the filters retain exactly the
original feasible terms, with their signs, proving
\eqref{eq:lowering-identity}.
\end{proof}

\begin{corollary}\label{cor:inflation-no-escape}
For every connected \(D\) and every integer \(t\ge0\), each
connected mate of \(I_tD\) has a connected deflation, and
\begin{equation}\label{eq:no-escape}
 I_{-t}\operatorname{Cl}(I_tD)\subseteq\operatorname{Cl}(D).
\end{equation}
Consequently an equality lost under inflation cannot return at a later
integer inflation.
\end{corollary}

\begin{proof}
The assertion is immediate at \(t=0\).  For \(t>0\), a connected
mate \(H\) has the same row and gap multisets as \(I_tD\).
Its rows exceed \(t\), and its gaps are at least \(t\).
Subtracting \(t\) preserves
\(0\le g_i-t<\min(r_i-t,r_{i+1}-t)\), so \(I_{-t}H\) is
connected.  The common row count is \(n\); apply
\Cref{lem:inflation-lowering} to both equal functions.
For permanence, an equality at a later parameter lowers to an equality
at every earlier nonnegative parameter.
\end{proof}

The inclusion in \eqref{eq:no-escape} is one-sided.  It does not assert
that inflation preserves every skew Schur equality.

\subsection{Zero-free classes and legal charts}

A connected zero-free nonline diagram has every row and every adjacent overlap
of length at least two.  Thus its erosion is connected, no row disappears,
and \(\mathcal I\mathcal E(S)=S\).  Connectedness of erosion alone would not
suffice for that last equality on a general unit-bearing diagram.

\begin{lemma}[Zero-free class transport]\label{lem:zero-free-bijection}
If \(S\) is connected, zero-free and not a line, erosion and inflation are
inverse bijections between \(\operatorname{Cl}(S)\) and
\(\operatorname{Cl}(\mathcal E(S))\).
\end{lemma}
\begin{proof}
Write a diagram as \(\lambda/\mu\).  The number of indices equal to
\(-1\) in its Jacobi--Trudi matrix is a Schur invariant on actual
connected diagrams.  For \(i=j+k-1\) such an entry says
\(\lambda_i-\mu_j=k-2\).  Diagonal and upper entries are positive,
and adjacent rows have positive overlap, so \(k=2\) cannot contribute.
For \(k\ge3\), this is a positive \(k\)-row overlap, whose multiplicity
is read from the overlap multiset \cite{Reiner2007Coincidences}.  Therefore
\[
 \#\{\text{index }-1\}
 =\sum_{k\ge3}\operatorname{mult}_{k-2}
   (\text{positive \(k\)-row overlaps}).
\]
The erosion of \(S\) has no \(-1\) entry, since \(S\) has no unit.  Neither
does any of its actual mates.  They also have the same row count, read
as the maximum ordinary degree.  On precisely this no-\(-1\) domain,
inflation is the whole homogeneous shift
\[
 s_{\mathcal I(T)}=h_1^{r(T)}s_T(h_j\mapsto h_{j+1}/h_1).
\]
Old units become \(h_1\), old negative indices are at most \(-2\) and remain
negative, and no new unit is omitted.  Conversely every mate of \(S\) is
zero-free, since a common beta would produce a lower ordinary-degree core
term, whereas \(s_S\) is homogeneous in ordinary degree.  Erosion and inflation
are inverse on the actual row profiles, and commute with rotation.
Together with \Cref{lem:erosion}, the raising identity gives both
directions on the entire classes, not just the one-sided inclusion
of \Cref{cor:inflation-no-escape}.
\end{proof}

\begin{lemma}[Inflation of legal charts]\label{lem:inflate-chart}
Suppose \(T=C\circ_\mu V\) is an actual proper chart, \(C\) is not a row, and
\(\mathcal I(T)\) is zero-free.  Write \(q,w\) for its period and width, with the
empty mark read horizontally.  Then \(\mathcal I(T)\) has an actual period-\(q\)
chart as follows:
\begin{enumerate}
\item for \(0<w<q-1\), use inner \(\mathcal I(V)\) and mark \(\mathcal I(W)\),
 with the same outer, displacement and attachment;
\item for \(w=0\), use inner \(\mathcal I(V)\) with its singleton \((a)\) end mark
 and the same outer;
\item for \(w=q-1\), the original inner is the \((H,H)\) one-middle-cell atom
 \(\mathcal A(W)\), and
 \(\mathcal I(T)=\mathcal I(C)\circ_{\varnothing}\mathcal I(W)\).
\end{enumerate}
\end{lemma}
\begin{proof}
Zero-freeness of \(\mathcal I(T)\) says that \(T\) has no index \(-1\), equivalently
\(B_T(j)A_T(j+1)=0\).  Inflation acts by
\(A_{\mathcal I(T)}(j)=A_T(j)\),
\(B_{\mathcal I(T)}(j)=B_T(j-1)\) and
\(H_{\mathcal I(T)}(j)=H_T(j)+B_T(j)\).
For \(0<w<q-1\) the required new private heights are
\(1+A_V(w+1)\) and \(1+B_V(q-1)\).  A vertical attachment or nonempty tail forces
the neighboring original bit, so the no-\(-1\) condition makes each of these
heights one.  The only exceptions are a horizontal lower attachment with
\(\sigma=w\) and a horizontal upper attachment with \(\tau=0\).  Their free interior
bits are \(B_C(j)\) and \(A_C(j)\); since a nonrow connected outer has a one in
each word, the same no-\(-1\) condition supplies the two missing pinches.
At \(w=0\), the joining slots read \(A_C(j),B_C(j)\).
Their nonzero occurrences force \(A_V(1)=B_V(q-1)=0\) by the
same no-\(-1\) condition.  The inflated inner therefore has horizontal
singleton marks at \(\operatorname{sw}(V)\) and
\(\operatorname{ne}(V)+(0,1)\), with the original displacement \(v_H\)
and new phases \((0,1)\).

At positive width, inflation commutes with the full intersection on the marked row range:
all those row intersections are nonempty, and their interval endpoints are
shifted by the same unit.  The new phases are
\((\tau,\sigma+1),(\tau+1,\sigma),(\tau+1,\sigma+1),(\tau,\sigma)\)
in cases \((a),(b),(c),(d)\).  In both positive and empty width,
the private pinches, thin marked tails and diagonal criterion prove
connected complements.  The hatted middle
remains two contents away from the periodic hatted mark.  Maximality follows
from the full intersection and the same extreme-displacement test.  Inflation
distributes over the complete copy union and, in the mixed cases, over every
supplemental copy.  This proves the first two assertions.

At \(w=q-1\) there is one private diagonal.  Unless both neighboring slots are
free, the no-\(-1\) condition either fails or forces an entire outer word to
zero, making \(C\) a row.  Thus \(\tau=0,\sigma=w\) and both attachments are horizontal.
The atom consists of \(W\), a translate of \(W\), and the single joining middle
cell.  If \(\mathcal A(W)\) denotes this atom, its displacement is
\(\operatorname{ne}(W)-\operatorname{sw}(W)+(0,2)\) and direct copy union gives
\[
 \mathcal I(C\circ_W\mathcal A(W))
   =\mathcal I(C)\circ_\varnothing\mathcal I(W).
\]
One can check the union by its copy origins; the row intervals
\(\mathcal E(A)\) and \(\mathcal E(A)+(0,1)\) cover \(A\) when every row survives.
The private singleton and empty middle hat prove all four legality clauses
for the atom.  The remaining no-\(-1\) conditions also make \(\mathcal I(C)\)
zero-free.
The inner \(\mathcal I(W)\) has span \(q\), and its horizontal
edge displacement is the displayed atom displacement, so the right-hand
side is a genuine legal empty-mark chart of period \(q\).
\end{proof}

\section{Rigidity of overlapped powers}\label{sec:overlap-powers}

Let \(\rho=(x_1,\ldots,x_k)\) be a positive composition with
\(k\ge2\), and let \(\delta\) be an integer satisfying
\(0<\delta<\min(x_1,x_k)\).  For integers \(m\ge1\), define
\[
 \rho^{\odot_\delta m}
 =(x_1,x_2,\ldots,x_{k-1},
   (x_1+x_k-\delta,x_2,\ldots,x_{k-1})^{m-1},x_k).
\]
Thus adjoining endpoint parts are merged with overlap \(\delta\),
not by the ordinary near-concatenation in \eqref{eq:ribbon-composition}.

\begin{lemma}\label{lem:overlap-power-rigidity}
For \(m\ge2\), the composition
\(\theta=\rho^{\odot_\delta m}\) has no factorization
\(\theta=\gamma\circ\beta\) with \(|\gamma|\ge2\) and
\(\ell(\beta)\ge2\).  Consequently
\[
 \mathcal R(\theta)=\{\theta,\theta^*\},
\]
with repetitions removed.
\end{lemma}

\begin{proof}
Put \(W=|\rho|\), \(Q=W-\delta\), and
\(T=|\theta|=mQ+\delta\).  If \(C\) is the proper cut set of
\(\rho\), the endpoint inequalities give
\begin{equation}\label{eq:overlap-cut-set}
 C\subseteq\{\delta+1,\ldots,Q-1\},\qquad
 C(\theta)=\bigcup_{j=0}^{m-1}(jQ+C).
\end{equation}
Let \(f\) be the cut indicator on the integers in \([0,T]\), with
\(f(0)=f(T)=0\).  It is \(Q\)-periodic wherever both positions are
in this interval, and
\begin{equation}\label{eq:initial-zero-interval}
 f(x)=0\qquad(0\le x\le\delta).
\end{equation}

Suppose \(\theta=\gamma\circ\beta\) as in the statement, and put
\(q=|\beta|\).  Then \(q\mid T\) and \(q\le T/2\).
Away from multiples of \(q\), the cut word repeats the internal cut
word of \(\beta\).  There are therefore
\(t_1,\ldots,t_{q-1}\in\{0,1\}\) such that
\[
 f(x)=t_{x\bmod q}\qquad(q\nmid x).
\]
If \(q\le\delta\), \eqref{eq:initial-zero-interval} makes every
\(t_r\) zero, contrary to \(\ell(\beta)\ge2\).  Hence
\(q>\delta\).

Set \(h=\gcd(Q,q)\).  Since \(q\mid mQ+\delta\),
\begin{equation}\label{eq:period-gcd}
 h\mid\delta,\qquad h\le\delta<q.
\end{equation}
Also \(Q<T/2\), so \(Q+q\le T\).  For each
\(1\le r<q\) with \(r+Q\not\equiv0\pmod q\), the two
available positions \(r,r+Q\) give
\begin{equation}\label{eq:cut-period-comparison}
 t_r=t_{r+Q\bmod q}.
\end{equation}
The cycles of addition by \(Q\) modulo \(q\) are the residue
classes modulo \(h\).  Removing the vertex zero leaves every other
cycle connected and the zero-residue cycle a connected path.
Thus \eqref{eq:cut-period-comparison} makes the indicators constant on
each such class, with zero omitted.  The initial zero positions
\(1,\ldots,h\) meet every class, including the zero class at
\(h<q\).  Every indicator is zero, again a contradiction.

By the ribbon factor-reversal classification, a nonrotation ribbon mate
would require a proper multipart right factor.  The obstruction just
proved excludes it.  Any remaining one-part factors only scale the parts
and introduce no additional reversal choice.
\end{proof}

\begin{lemma}\label{lem:merged-block-rigidity}
Under \eqref{eq:saturation-domain}, for every \(m\ge2\),
\[
 \operatorname{Cl}(A_m)=\{A_m,A_m^*\}.
\]
No conclusion about rigidity of \(A_1\) is asserted.
\end{lemma}

\begin{proof}
Every internal row of \(A_m\) exceeds \(2c\).  For a merged row
this follows from \(a+b-d>2c\); for an original internal row it is
one of \eqref{eq:saturation-domain}.  Hence \(A_m\) is Hessenberg.
Its signed maximal cofactor is \(h_c^{m(k-1)}\).
By \eqref{eq:hessenberg-membership} and the polynomial readings, any
connected mate is also Hessenberg, has \(m(k-1)+1\) rows, and has
constant gap \(c\), including when \(c=0\).

Subtract \(c\) from each row of \(A\), and put
\[
 \rho=(a-c,a_2-c,\ldots,a_{k-1}-c,b-c),\qquad \delta=d-c.
\]
Then the row word of \(A_m\), minus \(c\) in each position, is
\(\rho^{\odot_\delta m}\), and
\(0<\delta<\min(a-c,b-c)\).  Equality of the constant-gap functions
is therefore ribbon equality by \eqref{eq:constant-gap-substitution}.
\Cref{lem:overlap-power-rigidity} gives precisely the two displayed
orientations.
\end{proof}

\section{The repeated-block inverse}\label{sec:saturation-proof}

Under \eqref{eq:saturation-domain}, put
\(K=k-1\) and \(Q=|A|-Kc-d\).  The construction has
\begin{equation}\label{eq:saturation-readings}
 \begin{aligned}
 n&=MK+\ell,&
 |F_\alpha|&=M|A|-(M-\ell)d,\\
 N_{F_\alpha}&=MQ+d,&
 (-1)^{n-1}[h_{MQ+d}]s_{F_\alpha}
   &=h_c^{MK}h_d^{\ell-1}.
 \end{aligned}
\end{equation}
The strict endpoint inequalities give \(Q>0\).
All source rows exceed \(d\).  At an endpoint adjoining a larger seam,
its incident gaps have sum \(c+d\), and at every other internal row
they have sum \(2c\).  Thus the source is Hessenberg.

\begin{lemma}\label{lem:ribbon-endpoints}
Ribbon Schur equality preserves the unordered pair of endpoint parts.
Consequently every connected mate of \(A_j\), including \(j=1\),
has unordered endpoint-row pair \(\{a,b\}\).
\end{lemma}

\begin{proof}
In a composition factorization, the rightmost factor with at least two
parts supplies the endpoint parts.  Any one-part factors to its right
scale both.  Reversing factors preserves or exchanges this pair.
The ribbon classification therefore proves its invariance.

The maximal cofactor of \(A_j\) forces every mate to have constant
gap \(c\) and Hessenberg form, just as in
\Cref{lem:merged-block-rigidity}.  Apply
\eqref{eq:constant-gap-substitution} and add \(c\) back to the
endpoint parts.  This determines their unordered pair, not the interior
of a one-copy component.
\end{proof}

\begin{lemma}\label{lem:total-contraction-inverse}
Let \(H\) be any connected mate of \(F_\alpha\), with
\(M\ge2\).  After at most one global rotation, cutting its larger
seams gives the literal ordered components
\(A_{\beta_1},\ldots,A_{\beta_\ell}\), where \(\beta\) is a
permutation of the parts of \(\alpha\).
\end{lemma}

\begin{proof}
By \eqref{eq:saturation-readings} and the polynomial readings, \(H\)
is Hessenberg with \(MK\) gaps \(c\), \(\ell-1\) gaps \(d\),
and the source's row multiset.  When \(c=0\), the row count supplies
its multiplicity.  Every row therefore exceeds \(d\), as required by
\Cref{lem:seam-coefficients}.

Put \(B_j=s_{A_j}\).  Cutting all larger seams in both diagrams and
using \Cref{lem:seam-coefficients} gives
\begin{equation}\label{eq:cut-components}
 \prod_{i=1}^{\ell}s_{H_i}
 =[h_d^0]s_H
 =\prod_{i=1}^{\ell}B_{\alpha_i}.
\end{equation}
Connected irreducibility matches the component functions with
multiplicity.  A component with function \(B_j\) has \(jK+1\)
rows, identifying \(j\); in their actual order these indices form a
permutation \(\beta\) of the parts of \(\alpha\).
Every \(H_i\) has endpoint pair \(\{a,b\}\) by
\Cref{lem:ribbon-endpoints}, but a one-copy component may still be
any mate of \(A\).

Contract all larger seams of \(H\), obtaining \(C\).  The other
extreme coefficient is
\begin{equation}\label{eq:total-contraction}
 [h_d^{\ell-1}]s_H=(-1)^{\ell-1}s_C
 =(-1)^{\ell-1}B_M.
\end{equation}
By \Cref{lem:merged-block-rigidity}, \(C=A_M\) or \(A_M^*\).
Rotate \(H\) once if necessary, and henceforth fix the literal
ordered word
\[
 C=A_M=(a,I,(z,I)^{M-1},b;c^{MK}).
\]

For \(1\le j<\ell\), put \(s_j=\beta_1+\cdots+\beta_j\).
The \(j\)-th contracted seam occupies the row position
\begin{equation}\label{eq:split-positions}
 1+Ks_j
\end{equation}
in \(C\): the first \(j\) components had \(Ks_j+j\) rows, and
the \(j-1\) preceding contractions leave \(Ks_j+1\).
Each component has at least two rows, so distinct seams share no
endpoint row.  At every position \eqref{eq:split-positions}, the
contracted row has length \(z=a+b-d\).

Starting with the exterior row \(a\), the endpoint pair forces the
first component to end with \(b\), and
\[
 b+v-d=z=a+b-d\quad\Longrightarrow\quad v=a
\]
forces the next component to start with \(a\).  Induction fixes every
split, including when \(a=b\).  The unmerged rows retain their
positions in \(C=A_M\), so the pieces are exactly \(A_{\beta_i}\),
with every remaining gap equal to \(c\).

The reconstruction uses positions, not distinct row lengths: it allows
\(z\) or the endpoints to repeat among internal rows and does not
require one-copy rigidity.  For \(\ell=1\), there is no split and
\eqref{eq:total-contraction} gives the conclusion.
\end{proof}

\begin{proof}[Proof of \Cref{thm:repeated-block-saturation}]
By \Cref{lem:total-contraction-inverse}, an arbitrary mate is
\(F_\beta\) or \(F_\beta^*\), with the same multiset of outer
parts.  It remains to determine which orders are equal.
The largest generator in \(B_m=s_{A_m}\) is private and satisfies
\begin{equation}\label{eq:private-saturation-corner}
 B_m=(-1)^{mK}h_c^{mK}h_{mQ+d}+R_m,\qquad
 R_m\in\mathbb Q[h_1,\ldots,h_{mQ+d-1}].
\end{equation}
No \(B_m\) involves \(h_d\).  After localization at \(h_d\),
\Cref{lem:triangular-substitution} shows that
\(H_m\mapsto B_m/h_d\) is injective over \(\mathbb Q(h_d)\).
The seam expansion gives
\begin{equation}\label{eq:saturation-substitution}
 s_{F_\beta}=h_d^{\ell(\beta)}
 r_\beta(B_1/h_d,B_2/h_d,\ldots).
\end{equation}
Since the outer lengths agree, equality with the source is equivalent to
\(r_\beta=r_\alpha\).  Conversely every such ribbon equality
gives equality in \eqref{eq:saturation-substitution}, and rotation
preserves it.  This proves \eqref{eq:saturation-class}.

The gaps \(d>c\) uniquely parse \(F_\beta\) into blocks, and a
block's row count recovers its index.  Thus \(\beta\mapsto F_\beta\)
is injective.  If \(A=A^*\), then
\(F_\beta^*=F_{\beta^*}\), and the two sets in
\eqref{eq:saturation-class} coincide.  If an \(A_j\) were symmetric,
its endpoints would give \(a=b\), and its first and last internal
segments would give \(I=I^*\); hence \(A=A^*\).
When \(A\ne A^*\), parsing therefore shows that the two orientation
sets are disjoint.  This proves \eqref{eq:saturation-cardinality}.
The literal legal presentations are established next.
\end{proof}

\begin{proposition}\label{prop:saturation-chart}
Under \eqref{eq:saturation-domain}, the case-(a) chart with end-row
marks of length \(d\) satisfies Hypotheses I--V and gives
\(F_\alpha=\alpha\circ_{(d)}A\) literally, up to translation.
\end{proposition}

\begin{proof}
Normalize the southwest cell of \(A\) to \((k-1,0)\), put
\(S_A=|A|-(k-1)c\), and set \(T=S_A-k\).  The northeast cell
is \((0,T)\).  Use the marks
\begin{equation}\label{eq:saturation-marks}
 \begin{aligned}
 W_{\mathrm{sw}}&=\{(k-1,j):0\le j<d\},\\
 W_{\mathrm{ne}}&=\{(0,j):T-d+1\le j\le T\}.
 \end{aligned}
\end{equation}
Starting from a lower copy, the translations for an east step and a north
step of the outer ribbon are respectively
\begin{equation}\label{eq:saturation-translations}
 \Delta=(-k+1,T-d+1),\qquad
 \Delta-(1,1)=(-k,T-d).
\end{equation}
The first identifies the marks and merges two endpoint rows into
\(a+b-d\); the second leaves the rows separate with gap \(d\).
Following the ribbon row by row therefore gives exactly
\eqref{eq:repeated-blocks}.

The inequalities \(a-d>c\) and \(b-d>c\) put both marks in the
terminal free portions of their rows.  Removing either shortens that
row by \(d\), leaves its overlap with the adjacent row intact, and
preserves the weak orders of both row boundaries.  The complement is
therefore a connected skew diagram.  The marked content diagonals are
single-cell terminal diagonals, so neither placement admits a larger
copy on the same diagonals.  Their content intervals are
\([1-k,d-k]\) and \([T-d+1,T]\).  Since \(S_A>2d\), at least
one content diagonal separates them.

In the infinite amalgamated chain, the same endpoint inequalities leave
no northwest or southeast diagonal neighbor of a marked cell.  Thus
\(\overline W=\varnothing\), which gives Hypothesis IV.  Both
marks attach to the middle through one horizontal adjacency, giving
case (a) and Hypothesis V.  The original internal inequalities
\(a_i>2c\) keep the adjoining internal overlap strips apart throughout
the chain.  These observations verify I--V on the full parameter domain.
\end{proof}

\begin{remark}\label{rem:saturation-boundaries}
Both \(M\ge2\) and strict seam separation matter.  Consider
\begin{equation}\label{eq:ambiguous-inners}
 \begin{aligned}
 A^{(0)}&=(4,7,4,10,7;1^4),\\
 A^{(1)}&=(4,10,7,4,7;1^4).
 \end{aligned}
\end{equation}
After subtracting one from each row, these are three times the ribbon
compositions \((1,2)\circ(1,2)\) and \((2,1)\circ(1,2)\).
They are nonrotation-equal by factor reversal and
\eqref{eq:constant-gap-substitution}.  At \(M=1\), the proposed
conclusion would incorrectly assert rigidity of the inner itself.

At \(d=c=1\), the two parents
\(A^{(0)}\cdot_1 A^{(0)}\) and
\(A^{(1)}\cdot_1 A^{(1)}\) remain nonrotation-equal.  The same
substitution gives the ribbon compositions
\((1,1)\circ(3,6,3,9,6)\) and
\((1,1)\circ(3,9,6,3,6)\), still related by factor reversal.
At \(d=2\), both inners satisfy \eqref{eq:saturation-domain}; the
outer class of \((1,1)\) is a singleton, so each parent is
rotation-rigid.  The two parents are distinct and are not rotations,
hence their functions differ.  This does not assert that every equal-gap
context preserves ambiguity.
\end{remark}

\section{The thin two-row class}\label{sec:thin-and-lowering}

Throughout this section, \(1\le c<a<b\), and the profiles
\(E_m,G_\alpha\) are those of \eqref{eq:thin-domain}.  Their parameters
are
\begin{equation}\label{eq:thin-readings}
 n=M+\ell,\qquad |G_\alpha|=M(a+b),\qquad N_{G_\alpha}=Mq,
 \qquad
 (-1)^{n-1}[h_{Mq}]s_{G_\alpha}=h_c^M.
\end{equation}
In particular, the maximal cofactor suppresses exactly \(\ell-1\)
zero gaps, whose number is recovered from \(n\).

\begin{lemma}\label{lem:two-row-power-rigidity}
For every \(m\ge1\),
\[
 \operatorname{Cl}(E_m)=\{E_m,E_m^*\}.
\]
\end{lemma}

\begin{proof}
The profile \(E_m=(a,(a+b)^{m-1},b;c^m)\) is Hessenberg, since
\(a+b>2c\).  Its maximal cofactor is \(h_c^m\).
An arbitrary connected mate is therefore Hessenberg with the same row
multiset and constant gap \(c\).  The row of length \(a\) is its
unique shortest row.

The variable \(h_a\) occurs only at the corresponding diagonal
position.  It is not a gap variable, and every nontrivial interval label
exceeds the rows it contains.  In the source,
\begin{equation}\label{eq:short-row-coefficient}
 [h_a]s_{E_m}=s_{\overline E_m}.
\end{equation}
The connected tail makes this coefficient irreducible.  At an internal
\(a\)-row of a mate, deleting the corresponding row and column leaves
two nonempty connected diagonal blocks with zero lower-left block.
Its \(h_a\)-coefficient would therefore be reducible, a contradiction.
Rotate so that the \(a\)-row is first.  Equality of the full functions
and of their \(h_a\)-coefficients now gives both observations in
\eqref{eq:boundary-reconstruction-input}, which recovers \(E_m\)
literally.
\end{proof}

\begin{theorem}\label{thm:two-row-thin-class}
For every nonempty composition \(\alpha\),
\begin{equation}\label{eq:thin-class}
 \operatorname{Cl}(G_\alpha)
 =\{G_\beta,G_\beta^*: \beta\in\mathcal R(\alpha)\}.
\end{equation}
The two orientation sets are disjoint, and the class has
\(2|\mathcal R(\alpha)|\) members.
\end{theorem}

\begin{proof}
By \eqref{eq:thin-readings}, an arbitrary connected mate \(H\)
is Hessenberg with \(M\) gaps \(c\) and \(\ell-1\) zero gaps.
Cut the zero seams into \(H_1,\ldots,H_\ell\), and contract them
all to \(C\).  With \(B_m=s_{E_m}\), \eqref{eq:unit-extremes} gives
\[
 \begin{aligned}
 [u^0]\widehat s_H(u)&=\prod_i s_{H_i}
                     =\prod_i B_{\alpha_i},\\
 [u^{\ell-1}]\widehat s_H(u)&=(-1)^{\ell-1}s_C
                     =(-1)^{\ell-1}B_M.
 \end{aligned}
\]
Connected irreducibility and \Cref{lem:two-row-power-rigidity} give
\(H_i=E_{\beta_i}\) or \(E_{\beta_i}^*\), where the row counts
\(\beta_i+1\) identify a permutation \(\beta\) of the parts of
\(\alpha\).  The row multiset of \(C\) is
\(\{a,b\}\cup\{a+b\}^{M-1}\).
Call the two component orientations positive and negative, according
as their endpoint pair is \((a,b)\) or \((b,a)\).  Removing the common
\(M-\ell\) internal rows of length \(a+b\), the remaining multiset is
\begin{equation}\label{eq:orientation-multiset}
 \{a,b\}\cup\{a+b\}^{\ell-1}.
\end{equation}
A positive-to-negative transition would create a row \(2b>a+b\).
Without such a transition, mixed signs form a negative string followed
by a positive string: both exterior rows are \(b\), and every junction
is \(2a\) or \(a+b\).  Either case contradicts
\eqref{eq:orientation-multiset}, the latter because no row is \(a\).
Thus one global rotation gives \(H=G_\beta\).

The corner in \(B_m\) is
\[
 B_m=(-1)^m h_c^m h_{mq}+R_m,\qquad
 R_m\in\mathbb Q[h_1,\ldots,h_{mq-1}].
\]
Thus \(H_m\mapsto B_m\) is injective by
\Cref{lem:triangular-substitution}.  Repeated use of
\eqref{eq:cut-merge} at zero joins gives
\begin{equation}\label{eq:thin-substitution}
 s_{G_\beta}=r_\beta(B_1,B_2,\ldots).
\end{equation}
Equality is therefore equivalent to \(\beta\in\mathcal R(\alpha)\).
Zero seams uniquely parse the blocks, so this parametrization is
injective.  Every \(G_\beta\) starts with \(a\) and ends with
\(b\), while its rotation starts with \(b\).  Since \(a<b\), the
orientation sets are disjoint.
\end{proof}

The proof used zero seams and a private shortest row.  It did not set
\(d=0\) in \Cref{thm:repeated-block-saturation}, whose domain requires
\(d>c\).


\section{The exact inflation threshold}\label{sec:inflation-proof}

\begin{lemma}\label{lem:subthreshold-transport}
Under \eqref{eq:thin-domain}, for \(0\le t<p\),
\[
 \operatorname{Cl}(I_tG_\alpha)
 =I_t\bigl(\operatorname{Cl}(G_\alpha)\bigr).
\]
For \(0<t<p\), the literal chart
\(I_tG_\alpha=\alpha\circ_{(t)}E(t)\) satisfies Hypotheses I--V.
\end{lemma}

\begin{proof}
A second-subdiagonal index of \(G_\alpha\) has negative depth among
\[
 a-c=p,\qquad b-c,\qquad a+b-2c.
\]
The first two occur beside zero joins, and the third inside an \(E_m\).
Entries farther below the diagonal are strictly smaller.  Therefore no
negative entry opens for \(0\le t<p\).  When \(\ell\ge2\), the
first depth occurs at every first row following a zero join, so \(p\)
is exactly the first opening parameter.  For \(\ell=1\), that
particular depth need not occur.

For \(0<t<p\), retaining the homogeneous units gives
\begin{equation}\label{eq:subthreshold-shift}
 s_{I_tG_\alpha}
 =h_t^n s_{G_\alpha}
 \left(h_j=\frac{h_{j+t}}{h_t}\ (j\ge1)\right).
\end{equation}
All members of the class in \Cref{thm:two-row-thin-class} have the
same row count and satisfy the same negative-depth bound.  Thus every
original equality survives.  \Cref{cor:inflation-no-escape} gives the
reverse inclusion, proving the asserted set equality.  At \(t=0\)
it is the identity.

For the chart, realize \(E(t)=(a+t,b+t;c+t)\) as
\begin{equation}\label{eq:inflation-inner-cells}
 \{(0,j):b-c-1\le j\le q+t-2\}
 \cup\{(1,j):0\le j\le b+t-1\},
\end{equation}
with marks
\begin{equation}\label{eq:inflation-marks}
 \begin{aligned}
 W_{\mathrm{sw}}&=\{(1,j):0\le j<t\},\\
 W_{\mathrm{ne}}&=\{(0,j):q-1\le j\le q+t-2\}.
 \end{aligned}
\end{equation}
An east step uses translation \((-1,q-1)\), and a north step uses
\((-2,q-2)\).  Amalgamation merges endpoint rows into
\((a+t)+(b+t)-t=a+b+t\); a north step gives gap \(t\).
The resulting profile is exactly \(I_tG_\alpha\).

Since \(t<a-c<b-c\), both marks lie in terminal free portions.
Their deletion preserves the weak orders of the row boundaries and
leaves the overlap unchanged.  Their content diagonals are terminal
single-cell diagonals and have an intervening diagonal, proving I--III.
In the infinite amalgamated chain the marked cells have no northwest or
southeast diagonal neighbors, so \(\overline W=\varnothing\) and
IV holds.  Each mark attaches horizontally through one cell, giving
case (a) and V.  At \(t=0\), the horizontal empty-overlap convention
gives \(G_\alpha\) directly.  Together with the factor-reversal
classification recalled in \Cref{subsec:ribbon-substitution}, the class
identity already proved gives every subthreshold mate up to global rotation.
\end{proof}

For a composition \(\alpha\) of \(M\), define its proper suffix-sum
set
\begin{equation}\label{eq:suffix-set}
 T_\alpha=\{\alpha_\ell,\alpha_{\ell-1}+\alpha_\ell,
          \ldots,\alpha_2+\cdots+\alpha_\ell\}
 \subseteq\{1,\ldots,M-1\}.
\end{equation}
It has \(\ell-1\) members and determines \(\alpha\).  For
\(\ell=1\) it is empty.

\begin{lemma}\label{lem:boundary-core}
Let \(P_\alpha=s_{I_pG_\alpha}\).  Its lowest ordinary-degree
component has degree \(M+1\) and equals the function of a connected,
zero-free Hessenberg diagram \(C_\alpha=(U_\alpha;V_M)\), where
\begin{equation}\label{eq:boundary-core-lists}
 \begin{aligned}
 U_\alpha={}&\{qk:1\le k\le M,\ k\notin T_\alpha\}
       \cup\{qk+p:k\in T_\alpha\cup\{M\}\},\\
 V_M={}&\{qj-a:1\le j\le M\}\cup\{0\}.
 \end{aligned}
\end{equation}
Each union is sorted decreasingly.  In particular, the column-beta list
is independent of \(\alpha\).
\end{lemma}

\begin{proof}
The beta lists of a single block are
\[
 X(E_m)=\{mq\}\cup\{b+jq:0\le j<m\},\qquad
 Y(E_m)=\{b-c+jq:0\le j<m\}\cup\{0\}.
\]
In a zero join the preceding block's lists are translated by the span
of all subsequent blocks.  Thus
\begin{equation}\label{eq:thin-beta-lists}
 \begin{aligned}
 X(G_\alpha)&=\{b+jq:0\le j<M\}
       \cup\{qs:s\in T_\alpha\cup\{M\}\},\\
 Y(G_\alpha)&=\{b-c+jq:0\le j<M\}
       \cup\{qs:s\in T_\alpha\cup\{0\}\}.
 \end{aligned}
\end{equation}
At \(t=p=a-c\), the identity \(b+p=q\) gives
\begin{equation}\label{eq:boundary-beta-lists}
 \begin{aligned}
 X(I_pG_\alpha)&=\{q,2q,\ldots,Mq\}
       \cup\{qs+p:s\in T_\alpha\cup\{M\}\},\\
 Y(I_pG_\alpha)&=\{qj-a:1\le j\le M\}
       \cup\{qs:s\in T_\alpha\cup\{0\}\}.
 \end{aligned}
\end{equation}
The inequalities
\begin{equation}\label{eq:boundary-residue-separation}
 0<p<a<a+p<q
\end{equation}
imply that the common betas are exactly \(qT_\alpha\).
Indeed the row residues are \(0,p\), while the nonzero column
residue is \(-a\); neither \(p\) nor \(-a\) is zero, and they
are distinct modulo \(q\).

Delete every common beta from both lists.  The retained lists are
\eqref{eq:boundary-core-lists}.  Writing
\(\varepsilon_k=\mathbf1_{k\in T_\alpha}\), their order is
\begin{equation}\label{eq:boundary-core-ordered}
 \begin{aligned}
 U_\alpha&=(Mq+p,\ Mq,\ (M-1)q+p\varepsilon_{M-1},
                  \ldots,q+p\varepsilon_1),\\
 V_M&=(Mq-a,\ (M-1)q-a,\ldots,q-a,\ 0).
 \end{aligned}
\end{equation}
At \(M=1\), the row list is simply \((q+p,q)\).
The sorted diagonal indices are positive, and the subdiagonal indices
are \(a\) or \(a+p\).  All second-subdiagonal indices are at most
\(a+p-q=a-b<0\).  The lists are disjoint by
\eqref{eq:boundary-residue-separation}.  Thus the residual diagram is
connected, zero-free, and Hessenberg.

The parent has \(M+\ell\) rows and exactly \(\ell-1\) unit
entries.  For \(s\in T_\alpha\), let \(j\) be the number of
suffix sums larger than \(s\).  In \eqref{eq:boundary-beta-lists},
the row and column positions of \(qs\) are
\[
 M-s+j+3\quad\hbox{and}\quad M-s+j+1.
\]
Their difference is two.  Selecting all the unit entries therefore has
positive simultaneous deletion sign: the common entries have the same
relative order, and the sum of the row-column differences is even.
Its residual determinant is \(s_{C_\alpha}\), which is nonzero.
Every term selecting fewer units has larger ordinary degree.  Hence
\begin{equation}\label{eq:boundary-lowest-degree}
 P_\alpha^{[M+1]}=s_{C_\alpha},\qquad
 \min\deg_h P_\alpha=M+1,
\end{equation}
where the superscript denotes an ordinary homogeneous component.
\end{proof}

\begin{lemma}\label{lem:private-residue-column}
Fix a positive integer modulus \(q\), residue sets \(R,S\), and a
strict column-beta list \(V_1>\cdots>V_m=0\).  Suppose a connected
diagram \(C=(U;V)\) has all row-beta residues in \(R\), all
nonterminal column-beta residues in \(S\), and
\begin{equation}\label{eq:private-residue-condition}
 R\cap(R-S)=\varnothing\pmod q.
\end{equation}
Then the positive-index generators in \(s_C\) with residues in
\(R\) are exactly
\(h_{U_1},\ldots,h_{U_m}\).  Among diagrams with these fixed residue
constraints and column list, \(s_C\) therefore determines the entire
row-beta list.
\end{lemma}

\begin{proof}
The residue condition makes each \(h_{U_i}\) private to its
last-column position and excludes every other generator of residue
in \(R\).  Its coefficient is the signed cofactor obtained by
deleting row \(i\) and the last column.  That minor's sorted diagonal
has positive indices \(U_j-V_j\) for \(j<i\) and nonnegative
indices \(U_{j+1}-V_j\) for \(j\ge i\).
The strict sorted-diagonal reading used in
\eqref{eq:maximal-cofactor} therefore gives a unique minimum-square
monomial, so the coefficient is nonzero.  The support in residues
\(R\) consequently recovers every row beta.
\end{proof}

\begin{proof}[Proof of \Cref{thm:sharp-inflation}]
\Cref{thm:two-row-thin-class,lem:subthreshold-transport} prove the
first line of \eqref{eq:inflation-class}, including its geometric
presentations.  For \(C_\alpha\) at the boundary, take
\[
 R=\{0,p\},\qquad S=\{-a\},\qquad R-S=\{a,a+p\}.
\]
These residue sets are disjoint as required by
\eqref{eq:boundary-residue-separation}.  Therefore
\Cref{lem:private-residue-column} reads the suffix set directly:
\begin{equation}\label{eq:suffix-support-recovery}
 T_\alpha=
 \{k:1\le k<M,\ h_{qk+p}\text{ occurs in }s_{C_\alpha}\}.
\end{equation}
For fixed \(a,b,c,M\), equality of two boundary-core functions forces
equality of their suffix sets and hence of their compositions.

By \Cref{cor:inflation-no-escape}, an arbitrary connected mate of
\(I_pG_\alpha\), after at most one rotation, is \(I_pG_\beta\)
for some \(\beta\in\mathcal R(\alpha)\).
Equality of their lowest ordinary-degree components gives
\(s_{C_\alpha}=s_{C_\beta}\), so
\eqref{eq:suffix-support-recovery} forces \(\beta=\alpha\).
This proves the boundary rotation orbit.  Applying
\Cref{cor:inflation-no-escape} to \(I_pG_\alpha\) with increment
\(t-p\) extends the conclusion to \(t>p\), since inflation is
injective and commutes with rotation.

Parsing and the two different endpoint lengths give
\(2|\mathcal R(\alpha)|\) members below the threshold and exactly
two above it.  For \(\ell=1\), the original ribbon class is already
a singleton; the proof still applies, with empty suffix set and no
common-beta deletion.  Thus the theorem includes that case without
asserting a decrease in its class size.
\end{proof}

\begin{remark}\label{rem:chart-boundary}
The displayed subthreshold chart fails at \(t=p\).  Removing
\(W_{\mathrm{ne}}\) from \eqref{eq:inflation-inner-cells} leaves a
top row ending at \(q-2\), while the bottom row ends at
\(b+p-1=q-1\).  The complement is no longer a skew diagram.
This explains the geometric location of the boundary, but does not
prove rigidity.  The coefficient identity
\eqref{eq:boundary-lowest-degree} and support reconstruction
\eqref{eq:suffix-support-recovery} exclude arbitrary mates, including
ones that might admit a different chart.
\end{remark}

\section{Examples and verification scope}\label{sec:examples}

For saturation with \(c=1\), \(d=2\), and \(\alpha=(1,2)\),
the asymmetric inner \((4,5;1)\) gives four class members, while the
symmetric inner \((4,4;1)\) gives two.  The outer ribbon class is
\(\{(1,2),(2,1)\}\) in both cases.  Equal endpoints alone do not
imply the second count: an inner with \(a=b\) but \(I\ne I^*\)
still has the separate orientation copy.  Nor need the merger length
be new.  For example, \(A=(4,7,5;1,1)\), \(d=2\), has
\(z=7\) equal to its internal row length; the cut positions in
\eqref{eq:split-positions} still determine every split.

The nonrigid inners \eqref{eq:ambiguous-inners} illustrate a different
point.  Their equal-gap ambiguity and its separation at \(d=2\)
are consequences of the two arguments in
\Cref{rem:saturation-boundaries}.  An independent exact check enumerates
all \(28{,}350\) orders of the row and gap multisets of
\(A^{(0)}\cdot_2 A^{(0)}\) and finds precisely its two rotations.
This is a complete-class check at one nonrigid-inner instance, not a
proof for arbitrary inner order.  The same witness at \(d=3\)
would have \(a=c+d\), outside the strict endpoint hypothesis of
\Cref{thm:repeated-block-saturation}; no further instance is inferred
from it here.

For inflation, \(E=(2,3;1)\) has threshold one.  The first inflations
of \(G_{(1,2)}\) and \(G_{(2,1)}\) are
\begin{equation}\label{eq:first-boundary-example}
 \begin{aligned}
 P&=(3,4,3,6,4;2,1,2,2),\\
 Q&=(3,6,4,3,4;2,2,1,2).
 \end{aligned}
\end{equation}
Deleting their respective common betas gives
\[
 C_P=((13,12,9,4);(10,6,2,0)),\qquad
 C_Q=((13,12,8,5);(10,6,2,0)).
\]
The last-column variable \(h_9\) occurs in the first core with
coefficient \(-h_2(h_3h_6-h_2h_7)\), and does not occur in the
second.  The lowest ordinary-degree components already distinguish the
two parents.  By contrast, \(E=(3,4;1)\) has threshold two:
the four-member class for \(\alpha=(1,2)\) survives at both
\(t=0\) and \(t=1\), then its images split into two rotation
classes at \(t=2\).  The threshold is not universally one.

\begin{corollary}\label{cor:unbounded-class-splitting}
There are connected skew Schur classes of arbitrarily large cardinality
whose images under \(I_1\) split into rotation classes only.
\end{corollary}

\begin{proof}
Let \(\alpha_e\) be the composition of \(e\ge1\) copies of
\((1,2)\) under \(\circ\).  Classical ribbon factorization gives
\(|\mathcal R(\alpha_e)|=2^e\).
For \(E=(2,3;1)\), \Cref{thm:sharp-inflation} gives
\[
 |\operatorname{Cl}(G_{\alpha_e})|=2^{e+1},\qquad
 |\operatorname{Cl}(I_tG_{\alpha_e})|=2\quad(t\ge1).
\]
Injectivity of \(I_1\) leaves all \(2^{e+1}\) images distinct,
partitioned into \(2^e\) rotation classes.
\end{proof}

The theorems are proved for all parameters satisfying their hypotheses.
The paper-local script \texttt{check\_ref1\_class\_inflation\_exact.py}
checks the displayed examples directly from Jacobi--Trudi determinants.
For the complete-class checks, row and adjacent-overlap multisets bound
all possible connected mates.  Every ordering is considered; integer
determinant evaluations only reject unequal candidates, and every
remaining candidate is compared as an exact polynomial.  The checked
saturation cases include zero inner gaps, symmetric inners, a merger
length equal to an internal row, equal endpoints with a nonpalindromic
interior, and the nonrigid inner above.

The script also checks the boundary beta lists, common-beta deletion sign,
private-column support, and lowering on all compositions of total size
at most five for four specified parameter triples.  Complete-class
inflation checks cover the threshold-one and threshold-two examples.
For \(e=2\) in \Cref{cor:unbounded-class-splitting}, it verifies that
the eight constructed members become four distinct rotation pairs;
it does not independently enumerate all fourteen-row mates.  None of
these finite checks proves an unbounded assertion.

\section{Concluding remarks}\label{sec:conclusion}

Cutting a repeated-block diagram determines functions, while total
contraction and positional splitting can determine its literal blocks.
This distinction permits \Cref{thm:repeated-block-saturation} to retain
nonrigid inners.  Repetition and strict seam separation supply different
parts of that argument, as their respective boundary examples show.

Additive inflation has a different effect.  On the two-row family it
transports the entire class below the threshold and separates its
nonrotation members at the threshold.  It remains injective throughout,
so it does not universally define an operation on the quotient by skew
Schur equality.
The exact bijection in \Cref{lem:zero-free-bijection} applies to the
erosion class of a connected zero-free nonline diagram; it does not
remove this boundary for general Schur classes.

\begin{question}\label{prob:nonconstant-inners}
Which nonconstant-gap inner diagrams admit a repeated-block inverse
without an inner-rigidity hypothesis?  The present proof needs both a
rigid total contraction and a way to locate its cuts from the recovered
component sizes.
\end{question}

\begin{question}\label{prob:inflation-survivors}
Which repeated-block classes have a nontrivial proper subset of their
original nonrotation mates surviving an inflation boundary?  The two-row
theorem gives complete survival or rotation alone; its fixed-column
boundary inverse is not supplied here for arbitrary inners.
\end{question}

\section*{Declarations}

{\footnotesize
\noindent\textbf{Funding.} This research received no external funding.

\smallskip
\noindent\textbf{Competing interests.} The author has no relevant interests to disclose.

\smallskip
\noindent\textbf{Data availability.} No datasets were generated or analyzed.

\smallskip
\noindent\textbf{Computational verification.} The proofs do not depend on numerical computation.  Scripts used for verification are available at \url{https://github.com/johnnyshields/papers}.

\smallskip
\noindent\textbf{Generative AI use.} During the preparation of this work, the author used OpenAI's ChatGPT (GPT-6 Astra, GPT-5.6 Sol) and Anthropic's Claude (Opus 5.5, Opus 5) to assist with mathematical research, manuscript drafting, and the development of auxiliary verification code.  After using these tools, the author reviewed and edited as needed.  The author takes full responsibility for the content of the work.
\par}

\begin{thebibliography}{99}

\bibitem{Barekat2009CoincidencesCorrigendum}
F. Barekat, V. Reiner and S. van Willigenburg,
\emph{Corrigendum to ``Coincidences among skew Schur functions''},
Adv. Math. \textbf{220} (2009), no.~5, 1655--1656,
DOI: \href{https://doi.org/10.1016/j.aim.2008.10.010}{10.1016/\allowbreak j.aim.2008.10.010}.

\bibitem{Billera2006DecomposableCompositions}
L.~J. Billera, H.~Thomas and S.~van Willigenburg,
\emph{Decomposable compositions, symmetric quasisymmetric functions and equality of ribbon Schur functions},
Adv. Math. \textbf{204} (2006), no.~1, 204--240,
DOI: \href{https://doi.org/10.1016/j.aim.2005.05.014}{10.1016/\allowbreak j.aim.2005.05.014}.

\bibitem{GarciaGarcia2020ToeplitzMinors}
D.~Garc\'ia-Garc\'ia and M.~Tierz,
\emph{Toeplitz minors and specializations of skew Schur polynomials},
J. Combin. Theory Ser. A \textbf{172} (2020), 105201,
DOI: \href{https://doi.org/10.1016/j.jcta.2019.105201}{10.1016/\allowbreak j.jcta.2019.105201}.

\bibitem{McNamara2009Classification}
P.~R.~W. McNamara and S.~van Willigenburg,
\emph{Towards a combinatorial classification of skew Schur functions},
Trans. Amer. Math. Soc. \textbf{361} (2009), no.~8, 4437--4470,
DOI: \href{https://doi.org/10.1090/S0002-9947-09-04683-2}{10.1090/\allowbreak S0002-9947-09-04683-2}.

\bibitem{Reiner2007Coincidences}
V.~Reiner, K.~M. Shaw and S.~van Willigenburg,
\emph{Coincidences among skew Schur functions},
Adv. Math. \textbf{216} (2007), no.~1, 118--152,
DOI: \href{https://doi.org/10.1016/j.aim.2007.05.006}{10.1016/\allowbreak j.aim.2007.05.006}.

\bibitem{Shields2026SkewSchur}
J.~Shields,
\emph{Boundary reconstruction and cofactor rigidity for skew Schur functions},
preprint, 2026,
DOI: \href{https://doi.org/10.5281/zenodo.22723463}{10.5281/\allowbreak zenodo.22723463}.

\bibitem{Shields2026Cofactor}
J.~Shields,
\emph{Cofactor independence and dilation for skew Schur functions},
preprint, 2026,
DOI: \href{https://doi.org/10.5281/zenodo.22732784}{10.5281/\allowbreak zenodo.22732784}.

\bibitem{Shields2026Gluing}
J.~Shields,
\emph{Local common-factor gluing and reduction of composition necessity for skew Schur functions},
preprint, 2026,
DOI: \href{https://doi.org/10.5281/zenodo.23133702}{10.5281/\allowbreak zenodo.23133702}.

\bibitem{Shields2026Coarsenings}
J.~Shields,
\emph{Short coarsenings determine ribbon Schur equivalence},
preprint, 2026,
DOI: \href{https://doi.org/10.5281/zenodo.23644543}{10.5281/\allowbreak zenodo.\allowbreak 23644543}.

\bibitem{Shields2026NegativeCount}
J.~Shields,
\emph{Skew Schur equivalence by Jacobi--Trudi negative count},
preprint, 2026,
DOI: \href{https://doi.org/10.5281/zenodo.23133741}{10.5281/\allowbreak zenodo.23133741}.
\end{thebibliography}

\end{document}
